\textbf{主張} 素数は無限に存在する．

任意の自然数 $n$ に対して，$n$ より大きい素数があることを示せばよい．
（上限を決めても必ずその先に素数があれば，有限個では尽きない）

$M = n! + 1$ とおく．
$2 \le k \le n$ のどの $k$ も $n!$ を割り切るので，$M$ を $k$ で割ると 1 余る．

$M$ の素因数 $p$ を1つとる（$M \ge 2$ なので存在する）．
上より $p$ は $2, \dots, n$ のどれでもないので $p > n$．

よって $n$ より大きい素数が存在する．$\blacksquare$
